\Needspace{17\baselineskip}

\CNsection{8.1}{Why Random Access is Needed}

When a UE (User Equipment) powers on or transitions from idle mode, it faces fundamental challenges:

{\def\LTcaptype{none} % do not increment counter
\Needspace{9\baselineskip}
\begin{longtable}[c]{@{}
  >{\raggedright\arraybackslash\hspace{0pt}}m{(\linewidth - 2\tabcolsep) * \real{0.4583}}
  >{\raggedright\arraybackslash\hspace{0pt}}m{(\linewidth - 2\tabcolsep) * \real{0.5417}}@{}}
\toprule\noalign{}
\rowcolor{CNink}\CNth{Challenge} & \CNth{Description} \\
\CNheadrulesep
\endhead
\bottomrule\noalign{}
\endlastfoot
\textbf{No uplink timing alignment} & UE doesn’t know its propagation delay to the base station \\
\textbf{No allocated uplink resources} & UE has no scheduled grants for transmission \\
\textbf{No unique identity} & Network hasn’t assigned a C-RNTI yet \\
\end{longtable}
}

\CNsubsection{}{Scenarios Requiring Random Access}

\begin{enumerate}
\def\labelenumi{\arabic{enumi}.}
\tightlist
\item
  \textbf{Initial Access} — UE powering on, first connection to network
\item
  \textbf{Handover} — Moving to a new cell, need synchronization with target eNodeB (eNB)
\item
  \textbf{DL Data Arrival (Idle Mode)} — UE paged for incoming data, must re-establish UL sync
\item
  \textbf{UL Data Arrival} — UE in connected mode lost UL synchronization
\item
  \textbf{Positioning} — Timing measurements for location services
\end{enumerate}

\begin{CNquote}

\begin{CNinsight}

\textbf{Key Insight:} Random access is the ONLY mechanism for a UE to initiate uplink communication when it has no prior scheduling from the network.

\end{CNinsight}

\end{CNquote}

\Needspace{14\baselineskip}

\CNsection{8.2}{PRACH (Physical Random Access Channel)}

\CNchained\CNsubsection{}{Time-Frequency Location in Resource Grid}

The PRACH occupies a specific position in the LTE resource grid:

\begin{itemize}
\tightlist
\item
  \textbf{Frequency domain:} 6 consecutive PRBs (1.08~MHz bandwidth)
\item
  \textbf{Time domain:} Configurable subframes based on PRACH configuration index
\item
  \textbf{Location:} Specified by \CNcode{prach-FreqOffset} (PRB offset from lowest UL PRB)
\end{itemize}

\CNsubsection{}{PRACH Configuration Index}

The configuration index (0–63) determines:

\begin{itemize}
\tightlist
\item
  \textbf{Which subframes} carry PRACH opportunities
\item
  \textbf{Periodicity} of PRACH occasions (every 1, 2, 5, or 10~ms)
\item
  \textbf{Preamble format} used
\end{itemize}

{\def\LTcaptype{none} % do not increment counter
\Needspace{11\baselineskip}
\begin{longtable}[c]{@{}cccc@{}}
\toprule\noalign{}
\rowcolor{CNink}\CNth{Config Index} & \CNth{Preamble Format} & \CNth{Subframe(s)} & \CNth{Density} \\
\CNheadrulesep
\endhead
\bottomrule\noalign{}
\endlastfoot
0 & 0 & 1 & 1 per frame \\
3 & 0 & 1, 4, 7 & 3 per frame \\
15 & 0 & 0–9 & 10 per frame \\
48 & 4 (short) & 1 & 1 per frame \\
\end{longtable}
}

\Needspace{17\baselineskip}

\CNsubsection{}{Preamble Formats}

{\def\LTcaptype{none} % do not increment counter
\Needspace{13\baselineskip}
\begin{longtable}[c]{@{}
  >{\centering\arraybackslash\hspace{0pt}}m{(\linewidth - 8\tabcolsep) * \real{0.2000}}
  >{\centering\arraybackslash\hspace{0pt}}m{(\linewidth - 8\tabcolsep) * \real{0.2000}}
  >{\centering\arraybackslash\hspace{0pt}}m{(\linewidth - 8\tabcolsep) * \real{0.2000}}
  >{\centering\arraybackslash\hspace{0pt}}m{(\linewidth - 8\tabcolsep) * \real{0.2000}}
  >{\centering\arraybackslash\hspace{0pt}}m{(\linewidth - 8\tabcolsep) * \real{0.2000}}@{}}
\toprule\noalign{}
\rowcolor{CNink}\begin{minipage}[b]{\linewidth}\centering
\CNth{Format}
\end{minipage} & \begin{minipage}[b]{\linewidth}\centering
\CNth{Sequence Length (Ts)}
\end{minipage} & \begin{minipage}[b]{\linewidth}\centering
\CNth{CP Length (Ts)}
\end{minipage} & \begin{minipage}[b]{\linewidth}\centering
\CNth{Total Duration}
\end{minipage} & \begin{minipage}[b]{\linewidth}\centering
\CNth{Max Cell Radius}
\end{minipage} \\
\CNheadrulesep
\endhead
\bottomrule\noalign{}
\endlastfoot
0 & 24576 (0.8~ms) & 3168 (0.1~ms) & \textasciitilde{}1~ms & \textasciitilde{}14.5~km \\
1 & 24576 (0.8~ms) & 21024 (0.68~ms) & \textasciitilde{}2~ms & \textasciitilde{}77.3~km \\
2 & 2\ensuremath{\times}24576 (1.6~ms) & 6240 (0.2~ms) & \textasciitilde{}2~ms & \textasciitilde{}29.5~km \\
3 & 2\ensuremath{\times}24576 (1.6~ms) & 21024 (0.68~ms) & \textasciitilde{}3~ms & \textasciitilde{}100.2~km \\
4 (short) & 4096 (0.13~ms) & 448 (0.015~ms) & \textasciitilde{}0.15~ms & \textasciitilde{}1.4~km \\
\end{longtable}
}

\begin{CNquote}

\textbf{Design Trade-off:} Longer CP \ensuremath{\to} larger cell radius but fewer PRACH opportunities per frame.

\end{CNquote}

\Needspace{14\baselineskip}

\CNsection{8.3}{Preamble Sequences}

\CNchained\CNsubsection{}{Zadoff-Chu (ZC) Sequences}

The PRACH preamble uses Zadoff-Chu sequences defined as:

\CNdisplay{x_u(n) = e^{-j\pi u n(n+1)/N_{ZC}}}

where:

\begin{itemize}
\tightlist
\item
  $u$ = root sequence index ($1 \le u \le N_{\mathrm{ZC}} - 1$)
\item
  $N_{ZC}$ = sequence length (839 for long, 139 for short)
\end{itemize}

\Needspace{14\baselineskip}

\textbf{Why Zadoff-Chu?}

{\def\LTcaptype{none} % do not increment counter
\Needspace{11\baselineskip}
\begin{longtable}[c]{@{}
  >{\raggedright\arraybackslash\hspace{0pt}}m{(\linewidth - 2\tabcolsep) * \real{0.5263}}
  >{\raggedright\arraybackslash\hspace{0pt}}m{(\linewidth - 2\tabcolsep) * \real{0.4737}}@{}}
\toprule\noalign{}
\rowcolor{CNink}\CNth{Property} & \CNth{Benefit} \\
\CNheadrulesep
\endhead
\bottomrule\noalign{}
\endlastfoot
\textbf{Constant amplitude (zero PAPR)} & Maximizes PA efficiency at cell edge \\
\textbf{Ideal periodic autocorrelation} & Enables precise timing estimation \\
\textbf{Zero cross-correlation between roots} & Different root \ensuremath{\to} zero interference \\
\textbf{Flat frequency spectrum} & Robust in frequency-selective channels \\
\end{longtable}
}

\CNsubsection{64}{Preambles Per Cell}

Each cell has exactly \textbf{64 available preambles}, derived from:

\begin{enumerate}
\def\labelenumi{\arabic{enumi}.}
\tightlist
\item
  \textbf{Root Sequence Index} (\CNcode{rootSequenceIndex}, broadcast in SIB2)
\item
  \textbf{Cyclic Shifts} applied to root sequences
\end{enumerate}

\CNdisplay{\mathrm{Preamble}(v) = \mathrm{CyclicShift}\left(x_u,\; v \times N_{CS}\right)}

Where:

\begin{itemize}
\tightlist
\item
  \CNcode{N\_CS} = cyclic shift size (\CNcode{zeroCorrelationZoneConfig})
\item
  Number of preambles per root = $\lfloor N_{\mathrm{ZC}} / N_{\mathrm{CS}} \rfloor$
\item
  Multiple roots used if one root doesn’t provide 64 preambles
\end{itemize}

\CNsubsection{}{Cell Coverage and Cyclic Shift}

The cyclic shift \CNcode{N\_CS} determines the \textbf{zero-correlation zone}, which must exceed the maximum round-trip delay:

\CNdisplay{N_{CS} > \frac{2\, d_{\max}}{c\, T_{s,\mathrm{seq}}}}

where $d_{\max}$ is the maximum cell radius, $c$ the speed of light ($3 \times 10^{8}$ m/\allowbreak{}s) and $T_{s,\mathrm{seq}}$ the sequence sample duration.

\textbf{Example:} For \ensuremath{N_{\mathrm{CS}}} = 13 (format 0):

\begin{itemize}
\tightlist
\item
  Max distinguishable delay = $13 \times (1/1.25\,\mathrm{MHz}) = 10.4\ \mu\mathrm{s}$
\item
  Max cell radius = $10.4\ \mu\mathrm{s} \times c/2 \approx 1.56\ \mathrm{km}$
\item
  Preambles per root = $\lfloor 839/13 \rfloor = 64$ \ensuremath{\to} single root sufficient
\end{itemize}

\Needspace{14\baselineskip}

\CNkeepfig{m08-cbra}{8}

\CNsection{8.4}{Contention-Based Random Access (4-Step CBRA)}

\CNchained\CNsubsection{}{Procedure Overview}

\CNfigure{m08-cbra}{Four-step contention-based random access}

\Needspace{11\baselineskip}

\CNsubsection{}{Step-by-Step Details}

\CNchained\CNsubsubsection{Step 1 — Msg1: Preamble Transmission}

\begin{itemize}
\tightlist
\item
  UE selects a \textbf{random preamble} from the available set (contention preambles)
\item
  Transmits on PRACH at configured time-frequency resources
\item
  \textbf{RA-RNTI} derived from PRACH occasion: \CNcode{RA-RNTI =}\allowbreak{}\CNcode{ 1 + t\_}\allowbreak{}\CNcode{id + 10 × f\_}\allowbreak{}\CNcode{id}

  \begin{itemize}
  \tightlist
  \item
    \CNcode{t\_id} = subframe index (0–9)
  \item
    \CNcode{f\_id} = frequency index (0–5)
  \end{itemize}
\end{itemize}

\Needspace{16\baselineskip}

\CNsubsubsection{Step 2 — Msg2: Random Access Response (RAR)}

RAR contains (per responding preamble):

{\def\LTcaptype{none} % do not increment counter
\Needspace{11\baselineskip}
\begin{longtable}[c]{@{}
  >{\raggedright\arraybackslash\hspace{0pt}}m{(\linewidth - 4\tabcolsep) * \real{0.3182}}
  >{\centering\arraybackslash\hspace{0pt}}m{(\linewidth - 4\tabcolsep) * \real{0.2727}}
  >{\raggedright\arraybackslash\hspace{0pt}}m{(\linewidth - 4\tabcolsep) * \real{0.4091}}@{}}
\toprule\noalign{}
\rowcolor{CNink}\CNth{Field} & \begin{minipage}[b]{\linewidth}\centering
\CNth{Bits}
\end{minipage} & \CNth{Purpose} \\
\CNheadrulesep
\endhead
\bottomrule\noalign{}
\endlastfoot
Timing Advance & 11 & UL timing correction (0–1282, step = $16 \times T_s$) \\
UL Grant & 20 & Resources for Msg3 \\
TC-RNTI & 16 & Temporary identity \\
Backoff Indicator & 4 & Backoff if overloaded \\
\end{longtable}
}

\begin{itemize}
\tightlist
\item
  \textbf{RAR Window:} 2–10 subframes (configurable by \CNcode{ra-ResponseWindowSize})
\item
  If no RAR received \ensuremath{\to} back off and retry
\end{itemize}

\CNsubsubsection{Step 3 — Msg3: Scheduled Transmission}

\begin{itemize}
\tightlist
\item
  UE transmits on UL-SCH using the grant from Msg2
\item
  Contains \textbf{RRC Connection Request} with:

  \begin{itemize}
  \tightlist
  \item
    UE identity (S-TMSI or random value)
  \item
    Establishment cause
  \end{itemize}
\item
  HARQ enabled (max 5 retransmissions for Msg3)
\end{itemize}

\CNsubsubsection{Step 4 — Msg4: Contention Resolution}

\begin{itemize}
\tightlist
\item
  eNB echoes the UE identity received in Msg3
\item
  UE checks if the echoed identity matches its own:

  \begin{itemize}
  \tightlist
  \item
    \textbf{Match} \ensuremath{\to} RACH successful, TC-RNTI becomes C-RNTI
  \item
    \textbf{No match} \ensuremath{\to} Collision detected, restart from Step 1 with backoff
  \end{itemize}
\end{itemize}

\Needspace{14\baselineskip}

\Needspace{20\baselineskip}

\CNsection{8.5}{Contention-Free Random Access (CFRA)}

\CNchained\CNsubsection{}{When Used}

{\def\LTcaptype{none} % do not increment counter
\Needspace{11\baselineskip}
\begin{longtable}[c]{@{}
  >{\raggedright\arraybackslash\hspace{0pt}}m{(\linewidth - 2\tabcolsep) * \real{0.5556}}
  >{\raggedright\arraybackslash\hspace{0pt}}m{(\linewidth - 2\tabcolsep) * \real{0.4444}}@{}}
\toprule\noalign{}
\rowcolor{CNink}\CNth{Scenario} & \CNth{Reason} \\
\CNheadrulesep
\endhead
\bottomrule\noalign{}
\endlastfoot
\textbf{Handover} & Fast sync needed, can’t afford collision delay \\
\textbf{DL data arrival (connected)} & eNB knows UE, can pre-assign preamble \\
\textbf{Positioning (OTDOA)} & Precise timing required \\
\textbf{Beam failure recovery} & Quick recovery needed \\
\end{longtable}
}

\CNkeepfig{m08-cfra}{3}

\CNsubsection{}{Procedure}

\CNfigure{m08-cfra}{Contention-free random access}

\Needspace{18\baselineskip}

\CNsubsection{}{Key Differences from CBRA}

{\def\LTcaptype{none} % do not increment counter
\Needspace{14\baselineskip}
\begin{longtable}[c]{@{}lcc@{}}
\toprule\noalign{}
\rowcolor{CNink}\CNth{Aspect} & \CNth{Contention-Based} & \CNth{Contention-Free} \\
\CNheadrulesep
\endhead
\bottomrule\noalign{}
\endlastfoot
Steps & 4 & 2 \\
Preamble selection & Random & Assigned by eNB \\
Collision possible & Yes & No \\
Latency & \textasciitilde{}15~ms (success) & \textasciitilde{}5~ms \\
Use case & Initial access & Handover, DL data \\
Preamble pool & Shared (e.g., 54) & Dedicated (e.g., 10) \\
\end{longtable}
}

\Needspace{14\baselineskip}

\CNkeepfig{m08-timing-advance}{10}

\CNsection{8.6}{Timing Advance (TA)}

\CNchained\CNsubsection{}{Why Timing Advance is Needed}

In OFDMA uplink (SC-FDMA), all UE signals must arrive at the eNB \textbf{aligned within the cyclic prefix} to maintain orthogonality:

\CNfigure{m08-timing-advance}{Timing advance aligns the uplink signals of near and far UEs}

\CNsubsection{}{TA Mechanism}

\textbf{Initial TA (from RAR - Msg2):}

\begin{itemize}
\tightlist
\item
  11-bit value: 0–1282
\item
  Step size: $16 \times T_s = 16 \times (1/30.72\,\mathrm{MHz}) \approx$ \textbf{0.52~\ensuremath{\mu}s}
\item
  Represents one-way advance applied to UL transmission
\end{itemize}

\textbf{TA Maintenance (MAC CE - Timing Advance Command):}

\begin{itemize}
\tightlist
\item
  6-bit value: 0–63
\item
  Applied as: \CNcode{TA\_}\allowbreak{}\CNcode{new =}\allowbreak{}\CNcode{ TA\_}\allowbreak{}\CNcode{old + (TA\_}\allowbreak{}\CNcode{command − 31) × 16 × Ts}
\item
  Sent periodically via MAC Control Element
\item
  If no TA update received within \CNcode{timeAlignmentTimer} \ensuremath{\to} UL sync lost
\end{itemize}

\Needspace{25\baselineskip}

\CNsubsection{}{Cell Radius from Timing Advance}

The maximum TA corrects for the \textbf{round-trip propagation delay}:

\CNdisplay{
\begin{aligned}
\text{Round-trip delay} &= \mathrm{TA} \times 16 \times T_s\\
\text{One-way delay} &= \mathrm{TA} \times 16 \times T_s / 2\\
\text{Max distance} &= c \times \text{one-way delay}
\end{aligned}
}

\CNsubsubsection{Numerical Example: Maximum Cell Radius}

\textbf{For initial TA (11-bit, max = 1282):}\CNdisplay{
\begin{aligned}
\text{Max round-trip delay} &= 1282 \times 16 \times T_s = 1282 \times 16 \times \frac{1}{30.72 \times 10^{6}}\\
&= 1282 \times 0.5208\ \mu\text{s} = 667.7\ \mu\text{s}\\
\text{Max one-way delay} &= 667.7 / 2 = 333.8\ \mu\text{s}\\
\text{Max cell radius} &= 3 \times 10^{8} \times 333.8 \times 10^{-6} = 100.1\ \text{km}
\end{aligned}
}

\textbf{Verification with preamble format:}

{\def\LTcaptype{none} % do not increment counter
\Needspace{11\baselineskip}
\begin{longtable}[c]{@{}cccc@{}}
\toprule\noalign{}
\rowcolor{CNink}\CNth{Format} & \CNth{Guard Time (CP)} & \CNth{Max RTD} & \CNth{Max Cell Radius} \\
\CNheadrulesep
\endhead
\bottomrule\noalign{}
\endlastfoot
0 & 99.5~\ensuremath{\mu}s ($3168 \times T_s$) & 99.5~\ensuremath{\mu}s & \textbf{14.9~km} \\
1 & 684.4~\ensuremath{\mu}s ($21024 \times T_s$) & 684.4~\ensuremath{\mu}s & \textbf{102.6~km} \\
2 & 203.1~\ensuremath{\mu}s ($6240 \times T_s$) & 203.1~\ensuremath{\mu}s & \textbf{30.5~km} \\
3 & 684.4~\ensuremath{\mu}s ($21024 \times T_s$) & 684.4~\ensuremath{\mu}s & \textbf{102.6~km} \\
\end{longtable}
}

\begin{CNquote}

\textbf{Note:} The actual cell radius is limited by the \textbf{smaller} of: TA range OR preamble CP duration.

\end{CNquote}

\Needspace{14\baselineskip}

\CNsection{8.7}{Collisions and Backoff}

\CNchained\CNsubsection{}{Collision Probability}

When N UEs simultaneously attempt random access with 64 available preambles:

\textbf{Probability that a specific UE collides with at least one other UE:}

\CNdisplay{P_{collision} = 1 - \left(1 - \frac{1}{64}\right)^{N-1}}

\Needspace{19\baselineskip}

\CNsubsubsection{Numerical Examples}

{\def\LTcaptype{none} % do not increment counter
\Needspace{16\baselineskip}
\begin{longtable}[c]{@{}ccc@{}}
\toprule\noalign{}
\rowcolor{CNink}\CNth{N (simultaneous UEs)} & \CNth{P(collision) per UE} & \CNth{Expected colliding UEs} \\
\CNheadrulesep
\endhead
\bottomrule\noalign{}
\endlastfoot
2 & 1.56\% & 0.03 \\
5 & 6.10\% & 0.31 \\
10 & 13.0\% & 1.30 \\
20 & 25.5\% & 5.10 \\
30 & 36.8\% & 11.0 \\
50 & 54.0\% & 27.0 \\
64 & 63.2\% & 40.4 \\
\end{longtable}
}

\textbf{Example Calculation (N=10):}\CNdisplay{
\begin{aligned}
P(\text{collision}) &= 1 - \left(1 - \tfrac{1}{64}\right)^{10-1} = 1 - \left(\tfrac{63}{64}\right)^{9}\\
&= 1 - (0.98438)^{9} = 1 - 0.8681 = 0.1319 \approx 13.2\%
\end{aligned}
}

\CNkeepfig{m08-collision}{3}

\CNsubsection{}{Collision Scenario with Resolution}

\CNfigure{m08-collision}{Preamble collision and its resolution in Msg4}

\CNsubsection{}{Backoff Mechanism}

When a RACH attempt fails, the UE must wait before retrying:

\begin{enumerate}
\def\labelenumi{\arabic{enumi}.}
\tightlist
\item
  \textbf{Backoff Indicator (BI):} Broadcast in RAR subheader

  \begin{itemize}
  \tightlist
  \item
    Values: 0, 10, 20, 30, 40, 60, 80, 120, 160, 240, 320, 480, 960~ms
  \item
    UE selects random backoff: uniform [0, BI]
  \end{itemize}
\item
  \textbf{Power Ramping:} Each failed attempt increases transmit power

  \begin{itemize}
  \tightlist
  \item
    Step size: \CNcode{powerRampingStep} (2~dB typical)
  \item
    \CNcode{P\_}\allowbreak{}\CNcode{PRACH(attempt) =}\allowbreak{}\CNcode{ P\_}\allowbreak{}\CNcode{initial + (attempt − 1) × powerRampingStep}
  \item
    Maximum: \CNcode{preambleInitialReceivedTargetPower} + ramping
  \item
    Max attempts: \CNcode{preambleTransMax} (3, 4, 5, 6, 7, 8, 10, 20, 50, 100, 200)
  \end{itemize}
\item
  \textbf{Preamble re-selection:} New random preamble chosen for each attempt
\end{enumerate}

\CNkeepfig{m08-rach-attempts}{3}

\CNsubsection{}{RACH Attempt Timeline}

\CNfigure{m08-rach-attempts}{RACH attempts with power ramping and backoff}

\Needspace{14\baselineskip}

\Needspace{22\baselineskip}

\CNsection{8.8}{RACH Overload}

\CNchained\CNsubsection{}{Causes of RACH Overload}

{\def\LTcaptype{none} % do not increment counter
\Needspace{13\baselineskip}
\begin{longtable}[c]{@{}
  >{\raggedright\arraybackslash\hspace{0pt}}m{(\linewidth - 4\tabcolsep) * \real{0.2593}}
  >{\raggedright\arraybackslash\hspace{0pt}}m{(\linewidth - 4\tabcolsep) * \real{0.4815}}
  >{\raggedright\arraybackslash\hspace{0pt}}m{(\linewidth - 4\tabcolsep) * \real{0.2593}}@{}}
\toprule\noalign{}
\rowcolor{CNink}\CNth{Cause} & \CNth{Description} & \CNth{Scale} \\
\CNheadrulesep
\endhead
\bottomrule\noalign{}
\endlastfoot
\textbf{M2M/\allowbreak{}IoT} & Massive simultaneous device activation (e.g., smart meters reporting hourly) & 10,000–30,000 devices \\
\textbf{Paging storms} & Major event triggers mass DL data (emergency alert, popular content) & Thousands of UEs \\
\textbf{Disaster/\allowbreak{}emergency} & Everyone calls simultaneously after earthquake/\allowbreak{}event & Extreme burst \\
\textbf{Cell reselection} & Many UEs hand over simultaneously (train, stadium) & Hundreds \\
\textbf{Power outage recovery} & All devices reconnect simultaneously & Thousands \\
\end{longtable}
}

\CNsubsection{}{Overload Symptoms}

\begin{itemize}
\tightlist
\item
  High collision rate (\textgreater{}50\%)
\item
  Excessive preamble retransmissions \ensuremath{\to} increased interference
\item
  RAR capacity exhausted (limited DL resources for RAR)
\item
  Msg3/\allowbreak{}Msg4 resource starvation
\item
  Access delays from seconds to minutes
\end{itemize}

\Needspace{11\baselineskip}

\CNsubsection{}{Solutions}

\CNchained\CNsubsubsection{Access Class Barring (ACB)}

\begin{itemize}
\tightlist
\item
  UEs assigned to Access Classes 0–9 (plus special classes 11–15)
\item
  eNB broadcasts barring factor (0–1) and barring time
\item
  UE draws random number; if \textgreater{} barring factor \ensuremath{\to} barred for barring time
\item
  Spreads access attempts over time
\end{itemize}

\CNsubsubsection{Extended Access Barring (EAB)}

\begin{itemize}
\tightlist
\item
  For delay-tolerant M2M/\allowbreak{}IoT devices (AC 0–9)
\item
  eNB broadcasts EAB bitmap (one bit per AC)
\item
  Barred UEs wait for barring to be lifted
\item
  More aggressive than ACB for low-priority traffic
\end{itemize}

\CNsubsubsection{Back-off Tuning}

\begin{itemize}
\tightlist
\item
  Dynamically increase BI (up to 960~ms) during overload
\item
  Effectively rate-limits RACH attempts
\item
  Trade-off: individual access delay vs. system throughput
\end{itemize}

\Needspace{16\baselineskip}

\CNsubsubsection{Additional 3GPP Solutions}

{\def\LTcaptype{none} % do not increment counter
\Needspace{13\baselineskip}
\begin{longtable}[c]{@{}
  >{\centering\arraybackslash\hspace{0pt}}m{(\linewidth - 4\tabcolsep) * \real{0.2000}}
  >{\raggedright\arraybackslash\hspace{0pt}}m{(\linewidth - 4\tabcolsep) * \real{0.3600}}
  >{\raggedright\arraybackslash\hspace{0pt}}m{(\linewidth - 4\tabcolsep) * \real{0.4400}}@{}}
\toprule\noalign{}
\rowcolor{CNink}\begin{minipage}[b]{\linewidth}\centering
\CNth{Release}
\end{minipage} & \CNth{Solution} & \CNth{Mechanism} \\
\CNheadrulesep
\endhead
\bottomrule\noalign{}
\endlastfoot
R10 & ACB & Probabilistic barring \\
R11 & EAB & IoT/\allowbreak{}M2M specific barring \\
R12 & ACDC & App-level congestion control \\
R13 & Group paging & Reduce paging-triggered RACH \\
R14 & EDT (Early Data Transmission) & Data in Msg3 \ensuremath{\to} fewer full connections \\
\end{longtable}
}

\Needspace{16\baselineskip}

\CNkeepfig{m08-two-step}{12}

\CNsection{8.9}{5G NR Random Access}

\CNchained\CNsubsection{}{Key Enhancements over LTE}

\CNchained\CNsubsubsection{2-Step RACH (Rel-16): MsgA + MsgB}

Combines the 4-step procedure into 2 steps for \textbf{ultra-low latency}:

\CNfigure{m08-two-step}{Four-step LTE RACH versus two-step NR RACH}

\textbf{When to use 2-step vs 4-step:}

\begin{itemize}
\tightlist
\item
  2-step: Small cells, low load, URLLC requirements
\item
  4-step: Large cells (need TA before UL data), high load scenarios
\item
  Fallback: If MsgB not received \ensuremath{\to} fall back to 4-step
\end{itemize}

\CNkeepfig{m08-ssb-rach}{4}

\CNsubsubsection{Beam-Based RACH (SSB-to-RACH Mapping)}

In 5G NR with beamforming, the RACH procedure is \textbf{beam-aware}:

\CNfigure{m08-ssb-rach}{SSB-to-RACH mapping: one PRACH occasion per beam}

\begin{itemize}
\tightlist
\item
  UE selects \textbf{best SSB beam} (highest RSRP)
\item
  Transmits preamble on \textbf{associated RACH occasion}
\item
  gNB knows which beam direction the UE is located \ensuremath{\to} \textbf{beam correspondence}
\item
  Enables directional RAR transmission
\end{itemize}

\textbf{Configuration:}

\begin{itemize}
\tightlist
\item
  \CNcode{ssb-perRACH-OccasionAndCB-PreamblesPerSSB}: maps SSBs to RACH occasions
\item
  Multiple SSBs can share one RACH occasion, or one SSB can span multiple
\end{itemize}

\CNsubsubsection{Supplementary Uplink (SUL) RACH}

\begin{itemize}
\tightlist
\item
  Lower frequency band (e.g., 700~MHz) for extended coverage
\item
  UE selects SUL PRACH when \CNcode{rsrp-ThresholdSSB-SUL} condition met
\item
  Separate preamble format and configuration for SUL
\item
  Enables coverage extension without increasing NR UL power
\end{itemize}

\Needspace{24\baselineskip}

\CNsubsubsection{NR Preamble Formats}

{\def\LTcaptype{none} % do not increment counter
\Needspace{21\baselineskip}
\begin{longtable}[c]{@{}cccc@{}}
\toprule\noalign{}
\rowcolor{CNink}\CNth{Format} & \CNth{Sequence Length} & \CNth{Subcarrier Spacing} & \CNth{Use Case} \\
\CNheadrulesep
\endhead
\bottomrule\noalign{}
\endlastfoot
0 & 839 (long) & 1.25~kHz & Large cells (FR1) \\
1 & 839 (long) & 1.25~kHz & Very large cells \\
2 & 839 (long) & 1.25~kHz & Extreme coverage \\
3 & 839 (long) & 5~kHz & Medium cells \\
A1 & 139 (short) & 15/\allowbreak{}30~kHz & Small cells (FR1) \\
A2 & 139 (short) & 15/\allowbreak{}30~kHz & Normal cells (FR1) \\
A3 & 139 (short) & 15/\allowbreak{}30~kHz & Normal cells (FR1) \\
B4 & 139 (short) & 60/\allowbreak{}120~kHz & FR2 (mmWave) \\
C0 & 139 (short) & 15/\allowbreak{}30~kHz & Msg3 early TX \\
C2 & 139 (short) & 15/\allowbreak{}30~kHz & Msg3 early TX \\
\end{longtable}
}

\CNsubsubsection{Beam Failure Recovery (BFR)}

A special RACH-like procedure in NR:

\begin{enumerate}
\def\labelenumi{\arabic{enumi}.}
\tightlist
\item
  UE detects beam failure (RSRP below threshold for all configured beams)
\item
  UE selects a \textbf{new candidate beam} (from measured SSBs/\allowbreak{}CSI-RS)
\item
  UE transmits \textbf{dedicated CFRA preamble} on PRACH occasion of new beam
\item
  gNB responds with BFR response \ensuremath{\to} reconfigure beam pair
\end{enumerate}

\Needspace{14\baselineskip}

\Needspace{23\baselineskip}

\CNsection{8.10}{Summary and Key Equations}

\CNchained\CNsubsection{}{Critical Formulas}

{\def\LTcaptype{none} % do not increment counter
\Needspace{14\baselineskip}
\begin{longtable}[c]{@{}
  >{\raggedright\arraybackslash\hspace{0pt}}m{(\linewidth - 2\tabcolsep) * \real{0.5500}}
  >{\raggedright\arraybackslash\hspace{0pt}}m{(\linewidth - 2\tabcolsep) * \real{0.4500}}@{}}
\toprule\noalign{}
\rowcolor{CNink}\CNth{Parameter} & \CNth{Formula} \\
\CNheadrulesep
\endhead
\bottomrule\noalign{}
\endlastfoot
Collision probability & $P_c = 1 - (63/64)^{N-1}$ \\
RA-RNTI & $1 + t_{id} + 10 \times f_{id}$ \\
Timing Advance (initial) & $TA = N_{TA} \times 16 \times T_s$ \\
Max cell radius & $d_{max} = c \times N_{TA,max} \times 16 \times T_s / 2$ \\
Power ramping & $P_n = P_{init} + (n-1) \times \Delta_{ramp}$ \\
Preambles per root & $\lfloor N_{ZC} / N_{CS} \rfloor$ \\
\end{longtable}
}

\CNkeepfig{m08-tradeoffs}{3}

\CNsubsection{}{Design Trade-offs}

\CNfigure{m08-tradeoffs}{RACH design trade-offs}

\CNboxsection{Review Questions}

\begin{CNbox}[unbreakable]{q}{Review Questions}

\begin{enumerate}
\def\labelenumi{\arabic{enumi}.}
\tightlist
\item
  Why can’t the UE simply transmit on PUSCH without performing RACH first?
\item
  Calculate the collision probability when 25 UEs attempt simultaneous access.
\item
  What is the maximum cell radius for PRACH format 3? Show your calculation.
\item
  Explain why contention-free RACH is used during handover instead of contention-based.
\item
  A cell has \CNcode{zeroCorrelationZoneConfig} = 11 (\ensuremath{N_{\mathrm{CS}}} = 15). How many root sequences are needed to generate 64 preambles?
\item
  Describe the beam failure recovery procedure in 5G NR and its relation to RACH.
\item
  An IoT network has 10,000 devices that simultaneously attempt RACH. Describe two mechanisms to handle this overload.
\item
  Compare the latency of 2-step NR RACH vs. 4-step LTE RACH. When would you choose each?
\end{enumerate}

\end{CNbox}

\CNsection{}{References}

\begin{itemize}
\tightlist
\item
  3GPP TS 36.211 §5.7: Physical Random Access Channel (LTE)
\item
  3GPP TS 36.213 §6: Random Access Procedure (LTE)
\item
  3GPP TS 36.321 §5.1: Random Access Procedure (MAC)
\item
  3GPP TS 38.211 §6.3.3: PRACH (NR)
\item
  3GPP TS 38.213 §8: Random Access Procedure (NR)
\item
  3GPP TS 38.321 §5.1: Random Access Procedure (NR MAC)
\end{itemize}
